\section{Related Literature}\label{sec:literature}

The starting point is \citet{Economides1989}, who studies quality choice in a Hotelling duopoly and reports maximal variety differentiation together with minimal quality differentiation. In the three-stage location--quality--price game, the reported pure-strategy solution is obtained by solving the price stage on the regular shared-market branch, substituting those prices into quality-stage profits, and then solving the resulting first-order conditions. Our reassessment keeps that original noncooperative pricing game but evaluates quality choices against the full downstream price continuation, including exclusion branches.

The closest prior critique is \citet{Bunte1995} and the corresponding Chapter 3 discussion in \citet{Bunte1997}. This prior critique is more direct than a generic observation about Hotelling instability. Bunte explicitly states that the specifications and derivations correspond to those of \citet{Economides1989}, but argues that the maximum-differentiation conclusion disappears once a stability condition is imposed on the R\&D decision. Accordingly, the present paper makes no firstness claim about identifying a stability problem in Economides. Bunte's analysis, however, is organized around alternative resolutions of the downstream price-instability problem. His comparison distinguishes a Hotelling scenario in which firms recognize mutual dependence from a limit-price scenario; the ensuing variety and quality choices are derived under those pricing scenarios. The present paper asks a different question: if one keeps the original noncooperative price game itself, what is the global quality continuation at the maximal locations once the price game is re-solved after every quality deviation? The answer is the piecewise continuation in Lemma~\ref{lem:price}, the exact global-deviation threshold in Proposition~\ref{prop:correction}, and the additional exclusion equilibria in Proposition~\ref{prop:asymmetric}.

Bunte also cites \citet{Sorenson1995} for the R\&D stability condition and notes that Sorenson uses the same specification for the R\&D problem but does not analyze the variety decision. Sorenson studies product-improving R\&D prior to differentiated Bertrand competition and the persistence of market leadership. Thus Sorenson is directly relevant to the stability of the quality-investment problem, but it does not supply a location--quality--price reconstruction of the Economides three-stage result. The overlap therefore kills any claim that quality-stage stability had gone unnoticed; it does not by itself deliver the maximal-location piecewise price continuation, the $4/27$ threshold, or the two explicit exclusion quality equilibria established here.

A recent structural comparison is \citet{CohenHeifetz2024}. They study a covered spatial duopoly with quadratic transportation costs and exogenous vertical quality asymmetry, with firms choosing locations first and prices second. Their main object is how an exogenous quality gap shifts equilibrium locations and can eventually exclude the low-quality seller. The overlap with the present paper is therefore structural rather than exact: both analyses connect horizontal differentiation, vertical quality differences, and exclusion, and both engage the Economides line of research. The strategic object differs materially, however. Here quality is itself chosen in the intermediate stage of the original linear-transport three-stage game, and the issue is whether a finite endogenous quality deviation moves the downstream price continuation across an exclusion boundary. Cohen and Heifetz consequently do not supply the global quality-deviation test, the exact $4/27$ threshold, or the asymmetric quality equilibria proved here.

The broader horizontal--vertical differentiation literature provides the conceptual background. \citet{NevenThisse1990} allow firms to choose horizontal and vertical product characteristics before price competition and obtain max--min patterns of differentiation. \citet{Economides1993} studies quality variation in a circular model with entry and location decisions. \citet{DegryseIrmen2001} examine quality provision when quality changes perceived horizontal differentiation, while \citet{PigaPoyagoTheotoky2005} study a location--R\&D--price game with location-dependent R\&D spillovers. These contributions establish that quality and horizontal differentiation interact strategically. The narrower issue isolated here is different: when an upstream quality action can change which branch of the downstream price game is played, local optimization on a regular branch is not a global equilibrium check.

Finally, the price-stage discontinuity is connected to the classic instability of linear Hotelling price competition emphasized by \citet{DAspremontGabszewiczThisse1979}. We use that fact only to clarify the modern SPNE interpretation in Section~\ref{subsec:spne}. The headline results do not depend on resolving the full location game; they condition on the maximal locations reported by the original theorem and reassess the reduced quality game at those locations. The resulting contribution is therefore deliberately narrow: it does not replace the full literature on stability or product differentiation, but identifies what follows when the original noncooperative price continuation is treated globally rather than only on its regular branch.
